% ===== 论文精读笔记（AI 生成骨架，人工补充）=====
% 路径: ~/Projects/notes/src/papers/YYYY-MM-DD_作者_主题.tex
% 编译: latexmk -lualatex 本文件
\documentclass[11pt]{ctexart}
\usepackage[margin=2.5cm]{geometry}
\usepackage{paralist,amsmath,amssymb}
\usepackage[colorlinks=true]{hyperref}
\title{%%TITLE%%}
\author{精读笔记}
\date{\today}

\begin{document}
\maketitle

\begin{center}
\small
%%META%%  % 形如: arXiv:2401.xxxxx | 作者 | 年份 | 会议/期刊
\end{center}

\section{一句话总结}
% 领域 + 任务 + 方法类型 + 结论

\section{研究问题与动机}
\begin{compactitem}
  \item 
\end{compactitem}

\section{方法核心思想}
\begin{compactitem}
  \item 
\end{compactitem}

\section{实验与关键结果}
\begin{compactitem}
  \item 数据集：\quad 基线：
  \item 关键数字：
  \item 消融结论：
\end{compactitem}

\section{创新点}
\begin{compactitem}
  \item 
\end{compactitem}

\section{局限与未来工作}
\begin{compactitem}
  \item 
\end{compactitem}

\section{对我研究方向的启发}
\begin{compactitem}
  \item 
\end{compactitem}

\section{引用}
% \cite{key}（登记到 refs.bib）

\end{document}
